\documentclass{article}

\usepackage[dblblindworkshop, final]{neurips_2026}

\usepackage[utf8]{inputenc}
\usepackage[T1]{fontenc}
\usepackage{hyperref}
\usepackage{url}
\usepackage{booktabs}
\usepackage{amsmath}
\usepackage{graphicx}
\usepackage{amsfonts}
\usepackage{nicefrac}
\usepackage{microtype}

\hypersetup{
  pdftitle={Same Failure, Different Cause: Ungrounded Tool-Call Provenance in Simulated Agent Evaluation},
  pdfauthor={Abhimanyu Prasad}
}

\title{Same Failure, Different Cause: Ungrounded Tool-Call Provenance in Simulated Agent Evaluation}

\workshoptitle{Grounded User Simulation for Model Evaluation and Training (UserSim)}

\author{%
  Abhimanyu Prasad \\
  \texttt{abhiprd2000@gmail.com} \\
}

\begin{document}

\maketitle

\begin{abstract}
Tool-calling agents evaluated against LLM-simulated users frequently invoke tools with identifiers that no earlier tool result confirmed. We show that this single surface failure has different causes in different models, and that one of those causes would be concealed by a more faithful user simulator. On three matched $\tau^2$-bench retail tasks, with both agents facing the same simulated user, Qwen2.5-7B-Instruct leaves 93.3\% (14/15) of eligible identifiers unconfirmed, and every one of them had been stated by the simulated user. Llama-3.1-8B-Instruct does so less often (45.3\%, 43/95), and invents 48.8\% of them with no traceable source. Acting on a user-stated identifier without checking it is only wrong when the user is wrong. Our simulator often was: 88.4\% (38/43) of calls carrying such an identifier returned an error, which is why the failure is visible here at all. A simulator drawing identifiers from the task database would make those calls succeed, and an agent that never verifies would score like one that always does. Simulator choice also sets how much material exists for this failure to appear: holding the agent fixed, 59.3\% of the identifiers in Llama's tool calls had been stated by a 7B simulated user, against 14.6\% with a 1.5B one. The direction holds on all three matched tasks. We characterize two models over seven matched conversations and make no significance claim.
\end{abstract}

\section{Introduction}

Agentic benchmarks increasingly place a language model on both sides of the conversation: one model acts as the agent under test, another plays the user. This makes evaluation cheap and repeatable. It also makes the simulator part of the measuring instrument, and instruments have properties of their own.

Recent work establishes that the choice of user simulator moves the score. \citet{seshadri2026lost} run a user study across four countries against $\tau$-bench retail and find agent success varies by up to nine percentage points across user LLMs, with simulators serving as poorer proxies for some speaker populations than others. To isolate user variation, they hold the agent fixed. That design measures how much the simulator moves the score; because only one agent is present, it cannot ask whether the same score reflects different underlying failures in different agents.

We reverse the design. Two open-weight agents, one shared user simulator, three matched tasks. We do not report task success, because at this model scale there is none to report: six of the seven matched episodes scored a reward of $0.0$, and the seventh (Llama, task 3) could not be scored locally because its task requires $\tau^2$'s natural-language evaluator. Instead we characterize failure behavior directly, using a property that is measurable from transcripts regardless of whether the task was completed.

That property is grounding. When an agent calls a tool with an argument such as \texttt{order\_id}, did that value ever appear in the result of an earlier successful tool call in the same conversation? If not, the agent is acting on something it never confirmed.

Our contributions are:

\begin{enumerate}
\item Unconfirmed identifiers are common in both agents, but their \emph{provenance} differs by model. Every identifier Qwen used without confirmation, once confirmation was possible, had been stated by the user. Llama invents about half of its own.
\item One of these two failure modes would be concealed by a more faithful simulator. Our simulator stated incorrect identifiers often enough that the failure was observable here; one drawing identifiers from the task database would make every unverified claim true, and the agent that never checks would look identical to the one that always does.
\item We size the exposure this depends on. Holding the agent fixed, 59.3\% of the identifiers Llama acts on had been stated by a 7B simulator, against 14.6\% with a 1.5B one, so how often the failure can occur is partly a property of the simulator, not only of the agent.
\item We report a second, related failure axis: schema conformance. Llama malformed tool arguments in 10 of 10 opportunities and had two requests rejected outright by the parser; Qwen did neither in 8 opportunities. The malformation also interferes with provenance measurement, which we correct for (Section~\ref{sec:setup}).
\end{enumerate}

\section{Setup}
\label{sec:setup}

\paragraph{Harness.} $\tau^2$-bench v1.0.1 \citep{barres2025tau2}, retail domain. Every metric we report is computed from the conversation transcripts. The only evaluator output we cite is the task reward, and none of it derives from $\tau^2$'s natural-language assertions, which require a frontier API. Results from $\tau^2$-bench versions below 1.0.1 are not comparable with 1.0.1 and above; we pin the version and record it with every run.

\paragraph{Models.} Agents: Qwen2.5-7B-Instruct-AWQ and Llama-3.1-8B-Instruct-AWQ-INT4, both 4-bit. User simulator: Qwen2.5-7B-Instruct-AWQ for all matched runs. Section~\ref{sec:invisible} additionally uses three Llama conversations against a Qwen2.5-1.5B-Instruct simulator. Temperature $0$, \texttt{max\_steps} $200$.

\paragraph{Hardware.} Two NVIDIA T4 GPUs, 16\,GB each, one vLLM 0.27.1 server per GPU. All models are 4-bit quantized to fit within that budget, which also constrains the model sizes we can compare.

\paragraph{Data.} Seven matched conversations: four Qwen tasks (0--3) and three Llama tasks (0, 1, 3). Section~\ref{sec:results} uses the three tasks both agents completed. Llama task 2 could not be collected: the vLLM \texttt{llama3\_json} parser rejects the parallel tool call the model attempts, reproducibly, on two independent runs. We report this as a finding in Section~\ref{sec:schema} rather than working around it.

\paragraph{Grounding.} An identifier-shaped argument is \emph{grounded} if its exact string value appeared in the content of an earlier \emph{successful} tool result in the same conversation. Error results are excluded from the grounding pool: $\tau^2$'s error messages echo the offending identifier back, so including them would count a rejected value as evidence for itself.

We report two rates. The \emph{flat} rate is over all identifier arguments. The \emph{conditional} rate counts only calls made after at least one successful tool result already existed, so that the agent had something available to check against. The conditional rate is the honest one: a conversation in which no tool call ever succeeds has an empty grounding pool, and a 100\% ungrounded rate there is arithmetic rather than evidence of ignoring available information.

\paragraph{Provenance and terminology.} Each ungrounded identifier is assigned one source, checked in this order: the policy document, if the value appears there verbatim; otherwise a user utterance, if it appears verbatim in a user turn earlier than the call; otherwise \emph{invented}. Matching is exact, with no normalization. We call identifiers copied from the user \emph{unverified} and the rest \emph{invented}. Only the second is fabrication. Using a customer-supplied order number is reasonable; acting on it without a lookup is not verification. We hand-audited all 32 identifiers labeled invented against every preceding user turn, and all 32 labels held. Seventeen regression tests cover the grounding, provenance and parsing rules.

\paragraph{List-valued arguments.} Llama often emits a list-valued argument as a single string, such as \texttt{\textquotesingle[1151293680, 4983901480]\textquotesingle} (Section~\ref{sec:schema}). Treated as one opaque value, the bracketed string occurs nowhere else in the conversation, so it would always be scored as ungrounded and invented, even when every component had been returned by an earlier lookup or stated by the user. That would count one behavior, already scored as a schema error, a second time as fabrication. We therefore parse such strings into their components and classify each one. Qwen produced no stringified lists.

\paragraph{Denominator audit.} An initial heuristic identified arguments by an \texttt{\_id}/\texttt{\_ids} suffix. Auditing against the retail tool schema found that \texttt{email}, \texttt{zip}, \texttt{first\_name} and \texttt{last\_name} are also genuine identifier arguments. Including them raises the pooled conditional rate from 48.0\% to 51.8\% (both agents, matched tasks). All numbers below use the wider denominator.

\section{Results}
\label{sec:results}

\subsection{Grounding rate under a matched user simulator}

Both agents facing Qwen2.5-7B-Instruct-AWQ as the user, on tasks 0, 1 and 3.

\begin{table}[h]
\centering
\caption{Conditional ungrounded rate, matched user simulator, tasks 0, 1 and 3.}
\begin{tabular}{lcc}
\toprule
Agent & Conditional rate & Tasks with conditional data \\
\midrule
Qwen2.5-7B-Instruct-AWQ & 14/15 = 93.3\% & 2 of 3 \\
Llama-3.1-8B-Instruct-AWQ & 43/95 = 45.3\% & 3 of 3 \\
\bottomrule
\end{tabular}
\end{table}

The pooled gap is 48.1 percentage points. Per task:

\begin{table}[h]
\centering
\caption{Per-task conditional rates. Qwen's task 0 contributes nothing conditionally.}
\begin{tabular}{lccc}
\toprule
Task & Qwen & Llama & Gap \\
\midrule
0 & no data (pool never filled) & 22/50 = 44.0\% & --- \\
1 & 6/6 = 100.0\% & 15/35 = 42.9\% & 57.1 pp \\
3 & 8/9 = 88.9\% & 6/10 = 60.0\% & 28.9 pp \\
\bottomrule
\end{tabular}
\end{table}

Two caveats belong here rather than in the limitations section. Qwen's conditional evidence rests on two tasks, not three: in task 0 all eleven of its tool calls returned errors, so the grounding pool never became non-empty. And the gap is uneven across tasks: 57.1 points on task 1, 28.9 on task 3.

\paragraph{Statistics.} We ran an exact paired sign-flip permutation test over the three matched tasks. Because Qwen's task 0 has no conditional rate, the test necessarily uses the flat per-task rates and therefore \emph{does not test the headline conditional figure}. Per-task differences (Qwen minus Llama) are $+0.549$, $+0.556$, $+0.338$: all three positive. With three pairs there are exactly $2^3 = 8$ sign assignments, enumerated exhaustively. Only the observed assignment, with all three positive, is as extreme as the data, so the exact one-sided $p = 1/8 = 0.125$.

That value is the floor of this test. With three pairs the smallest attainable one-sided $p$ is $1/8$, so a perfectly consistent three-for-three result cannot reach conventional thresholds regardless of effect size. We report the finding as directionally consistent across all matched tasks and make no significance claim anywhere in this paper.

\subsection{Provenance}

Where did each ungrounded identifier come from? Table~\ref{tab:prov} gives both rates on the same matched tasks. The conditional rows correspond exactly to Table 1.

\begin{table}[h]
\centering
\caption{Provenance of ungrounded identifiers, matched tasks 0, 1 and 3. Neither model drew identifiers from the policy document.}
\label{tab:prov}
\begin{tabular}{llcc}
\toprule
Rate & Agent & From user (unverified) & Invented \\
\midrule
Flat & Qwen (54 ungrounded) & 43 (79.6\%) & 11 (20.4\%) \\
Flat & Llama (45 ungrounded) & 24 (53.3\%) & 21 (46.7\%) \\
Conditional & Qwen (14 ungrounded) & \textbf{14 (100.0\%)} & 0 (0.0\%) \\
Conditional & Llama (43 ungrounded) & 22 (51.2\%) & \textbf{21 (48.8\%)} \\
\bottomrule
\end{tabular}
\end{table}

This is the paper's central observation. The two models produce the same surface error at different rates, and for different reasons. Once any tool call has succeeded, Qwen never invents an identifier: all 14 of its conditional ungrounded identifiers were stated by the user, and all 11 of its invented ones occur before any tool call has succeeded. It takes what the customer said and acts on it without checking. Llama splits roughly evenly between the two.

These call for different fixes. ``Verify user-supplied values before acting'' and ``do not invent values'' are different instructions, and a benchmark reporting only an aggregate failure rate cannot indicate which one a given model needs.

\subsection{Why simulator fidelity decides which failures are observable}
\label{sec:invisible}

An agent that acts on a user-stated identifier without verifying it is harmed only when the user is wrong. Our simulator was often wrong. Across the matched tasks, 38 of the 43 tool calls carrying at least one unverified user-stated identifier returned an error (88.4\%), because the simulator stated identifiers absent from the task database. That is why this failure is visible in our data at all.

A simulator that stated only identifiers present in the database would invert the picture. Every unverified claim would happen to be true, every such call would succeed, and an agent that never verifies would be indistinguishable under evaluation from one that always does. The behavior would cost nothing on the benchmark and fail in deployment, where customers misremember order numbers. Simulator fidelity therefore determines not only how accurate a score is, but which failure modes are scoreable at all.

The opportunity for this failure also depends on the simulator. Table~\ref{tab:exposure} reports, for each conversation, the fraction of identifiers in the agent's tool calls that the simulated user had stated earlier. The first two rows hold the agent fixed and vary only the simulator.

\begin{table}[h]
\centering
\caption{Fraction of identifiers in the agent's tool calls that the simulated user had stated earlier in the conversation. $^\dagger$20-step smoke-test sample; excluding it, the 1.5B rate is 3/27 (11.1\%).}
\label{tab:exposure}
\begin{tabular}{llcccc}
\toprule
Agent & User simulator & Task 0 & Task 1 & Task 3 & Pooled \\
\midrule
Llama & Qwen2.5-7B & 9/10 & 5/10 & 2/7 & 16/27 (59.3\%) \\
Llama & Qwen2.5-1.5B & 3/14$^\dagger$ & 3/16 & 0/11 & 6/41 (14.6\%) \\
Qwen & Qwen2.5-7B & 5/11 & 8/9 & 4/7 & 17/27 (63.0\%) \\
\bottomrule
\end{tabular}
\end{table}

With the agent fixed, changing the simulator moves the share from 59.3\% to 14.6\%, and the one truncated run sits at the high end of the 1.5B range, so including it is conservative. A simulator that rarely states identifiers offers few opportunities for trust-without-verification to appear, whatever the agent's propensity. At the same 7B simulator, Qwen's share is 63.0\%, close to Llama's, so the simulator effect is much larger than the agent effect on this measure.

This is also why the matched design above was necessary. Our own initial comparison ran Qwen in self-play against the 7B user while Llama faced the 1.5B user. Under our classifier that comparison shows a 77.2-point gap; with a common simulator it is 48.1 points. Over a third of the original gap (37.7\%) was an artifact of simulator choice. The direction survived; the magnitude did not.

\subsection{Schema conformance: a second failure axis}
\label{sec:schema}

Grounding concerns \emph{which values} an agent uses. Conformance concerns \emph{whether the call is well-formed at all}. The two are distinct, but not independent: Section~\ref{sec:setup} shows how Llama's malformed arguments distorted its provenance numbers until they were parsed. The models separate on this axis too.

\begin{table}[h]
\centering
\caption{Schema conformance failures.}
\begin{tabular}{lccc}
\toprule
Agent & Malformed arguments & Type & Request-structure rejections \\
\midrule
Qwen & 0/8 & --- & 0 \\
Llama & \textbf{10/10} & string instead of list & 2 \\
\bottomrule
\end{tabular}
\end{table}

Llama's malformation is consistent: it emits \texttt{item\_ids:\ \textquotesingle[1151293680]\textquotesingle}, a string containing bracket characters, where the schema expects a list. Separately, on task 2 it attempts parallel tool calls, which the \texttt{llama3\_json} parser rejects with a hard error on both attempts. This is a model--parser interaction rather than a model-alone property, which means a deployment's parser choice partly determines the observed failure profile.

\subsection{Post-error behavior}

After a tool call fails, does the agent recover? Across 56 error events, the next action was a call to a genuine lookup tool in 29 (51.8\%), the same tool with different arguments in 10 (17.9\%), an exact repeat in 4 (7.1\%), something else in 9 (16.1\%), and nothing in 4 (7.1\%), where the episode ended. Calling a lookup tool is about twice as common as retrying (14, 25.0\%).

This does not repair the grounding rate. Even in conversations where lookup calls occur, subsequent calls remain largely ungrounded. Either the agent fails to extract the needed identifier from the lookup result, or the available lookup tools do not return the identifier type that is failing. Distinguishing these requires reading each lookup result against the next call's arguments, and we leave it to future work.

\section{Related Work}

\paragraph{Simulator robustness.} Our design responds directly to a limitation \citet{seshadri2026lost} name themselves: fixing the agent isolates user variation, but it forecloses any question about the agent. Their concern is representational, about which user populations a simulator stands in for. Ours is diagnostic, about which agent behaviors a simulator can expose at all. The two compound, since a simulator that misrepresents a population may also be blind to the failures that population would trigger.

\paragraph{Harness.} $\tau^2$-bench \citep{barres2025tau2} supplies our dual-control environment and its programmatic, database-state-based reward, extending the single-control setting of $\tau$-bench \citep{yao2024tau}.

\paragraph{Simulator fidelity.} \citet{chen2026simulated} measures decision fidelity against real purchase outcomes and finds simulated customers do not disengage the way real ones do. \citet{nakash2026efficient} identify a related over-cooperativeness in simulators. Our argument has that shape: a simulator that stated only correct identifiers would be the well-behaved one, and its good behavior is what would conceal the failure. Where those papers show that simulators are unfaithful, we show that one specific fidelity gap renders a whole class of agent failure unmeasurable.

\paragraph{Positioning.} Prior work varies the simulator against a fixed agent and measures how the score moves. We hold the simulator fixed and vary the agent. That reversal surfaces something the standard design cannot: the same surface failure has different provenance in different models, and one provenance is structurally invisible to any simulator that states identifiers correctly.

\section{Limitations}

\begin{enumerate}
\item \textbf{Two models, three matched tasks, seven conversations}, plus three Llama conversations against a 1.5B simulator. This is a characterization, not an estimate of a population rate. Confidence intervals would be meaningless at this scale and we report none.
\item \textbf{Qwen's conditional evidence comes from two of three matched tasks}, and its 15 conditional observations come from just two conversations (6 and 9).
\item \textbf{No significance claim is available.} The exact permutation floor at three pairs is $p = 1/8 = 0.125$.
\item \textbf{The significance test does not test the headline number.} It runs on flat rates, because Qwen's task 0 has no conditional rate.
\item \textbf{The grounding rules are hand-built.} Two implementation bugs in them were found by hand-tracing transcripts and corrected: one excluded some successful tool results from the grounding pool, the other missed legitimate reuse of identifiers containing a \texttt{\#}. Both inflated earlier rates. The rules, including list parsing, are covered by regression tests.
\item \textbf{The same model plays the user and one agent.} Qwen2.5-7B is both the simulator and the Qwen agent, so Qwen's tendency to copy user-stated identifiers may partly reflect affinity between outputs of the same model. A simulator from a different model family would test this.
\item \textbf{This analysis was not pre-registered.} The repository contains a plan pre-registered for an earlier rank-inversion design, abandoned when agents could not complete tasks at this scale, and an amendment recording why. The grounding analysis reported here was defined afterwards, from the same transcripts, and is exploratory.
\item \textbf{One sample per condition at temperature 0.} We did not vary the user's phrasing, so we cannot say how stable the provenance pattern is across rewordings of the same task.
\item \textbf{4-bit AWQ quantization on T4 GPUs.} We have no fp16 baseline and cannot rule out that quantization contributes to either failure mode.
\item \textbf{No matched episode was scored as a success.} Six scored $0.0$; Llama's task-3 episode could not be scored locally. This paper characterizes failure behavior, not task performance, and the findings may not transfer to models that complete tasks.
\item \textbf{Llama task 2 is missing} because of a parser rejection, so the matched comparison covers three tasks rather than four.
\item \textbf{The invisibility argument in Section~\ref{sec:invisible} is a counterfactual, not an observation.} Our simulator stated incorrect identifiers often enough that the trust-without-verification failure was visible here. We argue it would be concealed by a simulator whose stated identifiers are drawn from the task database, but we did not build such a simulator and did not measure that condition. Doing so is the most direct next experiment.
\end{enumerate}

\section{Reproducibility}

All runs used open-weight 4-bit models on two T4 GPUs. All reported metrics are computed locally from transcripts, with no external API dependency. $\tau^2$'s natural-language assertion evaluator does require a frontier API; we report nothing derived from it. Code, the ten raw transcripts, run records and seventeen regression tests are available at \url{https://github.com/abhiprd2000/simswap}. Every number in this paper is produced by scripts in \texttt{src/}; \texttt{docs/CAMERA\_READY\_NUMBERS.md} maps each one to its source.


\begin{thebibliography}{99}

\bibitem[Barres et al.(2025)]{barres2025tau2}
Victor Barres, Honghua Dong, Soham Ray, Xujie Si, and Karthik Narasimhan.
\newblock $\tau^2$-Bench: Evaluating conversational agents in a dual-control environment.
\newblock \emph{arXiv preprint arXiv:2506.07982}, 2025.

\bibitem[Chen(2026)]{chen2026simulated}
Liang Chen.
\newblock Simulated customers never walk away: Decision fidelity of LLM user
  simulators measured against real purchase outcomes.
\newblock \emph{arXiv preprint arXiv:2606.20708}, 2026.

\bibitem[Nakash et al.(2026)]{nakash2026efficient}
Itay Nakash, George Kour, and Ateret Anaby-Tavor.
\newblock Efficient agent evaluation via diversity-guided user simulation.
\newblock \emph{arXiv preprint arXiv:2604.21480}, 2026.

\bibitem[Seshadri et al.(2026)]{seshadri2026lost}
Preethi Seshadri, Samuel Cahyawijaya, Ayomide Odumakinde, Sameer Singh, and Seraphina Goldfarb-Tarrant.
\newblock Lost in simulation: LLM-simulated users are unreliable proxies for human users in agentic evaluations.
\newblock In \emph{Proceedings of the 64th Annual Meeting of the Association for Computational Linguistics}, 2026.

\bibitem[Yao et al.(2024)]{yao2024tau}
Shunyu Yao, Noah Shinn, Pedram Razavi, and Karthik Narasimhan.
\newblock $\tau$-bench: A benchmark for tool-agent-user interaction in real-world domains.
\newblock \emph{arXiv preprint arXiv:2406.12045}, 2024.

\end{thebibliography}

\end{document}